\documentclass[pdf, unicode, 12pt, a4paper, oneside, fleqn]{article}

\usepackage[utf8]{inputenc}
\usepackage[T2A]{fontenc}
\usepackage[english,russian]{babel}
\usepackage{amsmath}
\usepackage{amssymb}
\usepackage{array}
\usepackage{csquotes}
\usepackage{float}
\usepackage{graphicx}
\usepackage{hyperref}
\usepackage{longtable}
\usepackage{xcolor}

\frenchspacing
\setcounter{secnumdepth}{0}

\oddsidemargin=-0.4mm
\textwidth=160mm
\topmargin=4.6mm
\textheight=230mm

\parindent=0pt
\parskip=5pt

\newcommand{\code}[1]{\texttt{\detokenize{#1}}}
\newcommand{\placeholderfigure}[2]{
\begin{figure}[H]
    \centering
    \fbox{
        \parbox[c][0.27\textheight][c]{0.9\linewidth}{
            \centering
            #2
        }
    }
    \caption{#1}
\end{figure}
}

\begin{document}

\begin{titlepage}
\begin{center}
\bfseries

{\Large Московский авиационный институт\\
(национальный исследовательский университет)}

\vspace{48pt}

{\large Факультет информационных технологий и прикладной математики}

\vspace{36pt}

{\large Кафедра вычислительной математики и программирования}

\vspace{48pt}

Лабораторная работа по курсу \enquote{Информационный поиск}

\vspace{24pt}

{\Large Лабораторная работа \textnumero 1}\\
\enquote{Добыча корпуса документов}

\end{center}

\vspace{72pt}

\begin{flushright}
\begin{tabular}{rl}
Студент: & К.\,А. Арусланов \\
Преподаватель: & А.\,А. Кухтичев \\
Группа: & \underline{\hspace{55mm}} \\
Дата: & \underline{\hspace{55mm}} \\
Оценка: & \underline{\hspace{55mm}} \\
Подпись: & \underline{\hspace{55mm}} \\
\end{tabular}
\end{flushright}

\vfill

\begin{center}
\bfseries
Москва, \the\year
\end{center}
\end{titlepage}

\tableofcontents
\pagebreak

\section{Постановка задачи}

Необходимо подготовить корпус документов, который будет использован в следующих лабораторных работах по информационному поиску. Для выбранного набора документов требуется:

\begin{itemize}
    \item скачать примеры документов на локальный компьютер;
    \item указать источники данных;
    \item изучить структуру документов, разметку и доступную метаинформацию;
    \item выделить полезный текст из сырого HTML;
    \item найти существующие поисковые системы, которые уже позволяют искать по выбранным документам;
    \item привести примеры запросов и отметить недостатки существующей поисковой выдачи;
    \item посчитать статистику корпуса: размер сырых документов, количество документов, размер выделенного текста, средний размер документа и средний объём текста.
\end{itemize}

\section{Описание корпуса}

В качестве темы корпуса выбраны тексты песен. Один документ соответствует одной странице с текстом одной песни.

Для диагностической выборки были использованы три независимых источника:

\begin{itemize}
    \item Genius: \href{https://genius.com}{genius.com};
    \item AZLyrics: \href{https://www.azlyrics.com}{azlyrics.com};
    \item Musixmatch: \href{https://www.musixmatch.com}{musixmatch.com}.
\end{itemize}

На этапе лабораторной работы собиралась небольшая проверочная выборка: по 4 страницы из каждого источника. Такой объём достаточен для анализа формата документов, проверки парсеров и получения первичной статистики.

\section{Сбор данных}

Список страниц хранится в файле \code{urls.txt}. Формат каждой непустой строки:

\begin{verbatim}
source<TAB>url
\end{verbatim}

Например:

\begin{verbatim}
genius    https://genius.com/...
azlyrics  https://www.azlyrics.com/lyrics/.../....html
musixmatch        https://www.musixmatch.com/lyrics/...
\end{verbatim}

Сбор выполняется командой:

\begin{verbatim}
python main.py collect urls.txt
\end{verbatim}

Скрипт сохраняет сырой HTML в каталог \code{data/raw/<source>/}, а результат разбора в JSON-файлы каталога \code{data/parsed/<source>/}. Для каждого документа сохраняются URL, источник, HTTP-статус, размер сырого файла, извлечённый текст, размер текста, метаданные и SHA-256 выделенного текста.

Статистика строится командой:

\begin{verbatim}
python main.py stats
\end{verbatim}

Результаты сохраняются в файлы \code{results/statistics.json} и \code{results/statistics.csv}.

\section{Структура документов}

\subsection{Genius}

В сыром HTML Genius полезный текст находится в блоках:

\begin{verbatim}
<div data-lyrics-container="true">...</div>
\end{verbatim}

На странице также присутствуют служебные элементы интерфейса, например блоки с авторами разметки, переводами и описанием песни. Такие фрагменты помечаются атрибутом \code{data-exclude-from-selection="true"} и удаляются перед извлечением текста.

Метаинформация доступна в нескольких местах:

\begin{itemize}
    \item \code{meta[property="og:title"]}: содержит строку вида \code{Artist - Title};
    \item структурированные данные JSON-LD, если они присутствуют на странице;
    \item элементы шапки песни: название, исполнитель, дата релиза, альбом, просмотры;
    \item служебные данные приложения Genius, где встречаются теги жанров, аннотации, вопросы и дополнительные сведения о песне.
\end{itemize}

\subsection{AZLyrics}

AZLyrics имеет более простую HTML-структуру. Текст песни расположен в безымянном \code{div}, который следует после характерного HTML-комментария:

\begin{verbatim}
<!-- Usage of azlyrics.com content ... -->
\end{verbatim}

Название песни извлекается из \code{h1}, исполнитель --- из заголовка \code{h2}. Информация об альбоме может находиться рядом с текстом в служебной подписи страницы.

Особенность источника: структура страницы достаточно статична, но в HTML почти нет семантических классов для текста песни, поэтому парсер опирается на положение блока относительно комментария.

\subsection{Musixmatch}

Musixmatch использует Next.js. Визуально текст может отображаться в сгенерированных \code{div}-элементах с классами вида \code{css-g5y9jx}, однако такие классы нестабильны и не должны использоваться как основной признак.

Надёжнее извлекать данные из JSON-блока:

\begin{verbatim}
<script id="__NEXT_DATA__" type="application/json">...</script>
\end{verbatim}

Внутри этого объекта полезные данные находятся в ветке \code{props.pageProps.data.trackInfo.data}:

\begin{itemize}
    \item \code{lyrics.body}: полный текст песни;
    \item \code{track.name}: название трека;
    \item \code{track.artistName}: исполнитель;
    \item \code{track.albumName}: альбом;
    \item \code{track.releaseDate}: дата релиза;
    \item \code{trackStructureList}: структура текста по секциям и строкам;
    \item \code{lyrics.language}, \code{lyrics.verifier}, \code{lyrics.copyright}: дополнительные метаданные.
\end{itemize}

\section{Выделение текста}

Все источники приводятся к единому JSON-представлению:

\begin{verbatim}
{
  "source": "...",
  "url": "...",
  "title": "...",
  "artist": "...",
  "album": "...",
  "release_date": "...",
  "lyrics": "...",
  "metadata": {},
  "parse_error": null,
  "http_status": 200,
  "raw_size_bytes": 123456,
  "text_size_bytes": 4567,
  "lyrics_sha256": "..."
}
\end{verbatim}

При очистке текста выполняются следующие действия:

\begin{itemize}
    \item HTML-теги \code{br} заменяются на переводы строк;
    \item лишние пробелы и пустые строки нормализуются;
    \item служебные блоки интерфейса удаляются;
    \item отсутствующие метаданные не придумываются, а остаются пустыми.
\end{itemize}

\section{Статистика корпуса}

В таблице приведены результаты для диагностической выборки из 12 документов.

\begin{center}
\begin{tabular}{|l|r|}
\hline
Показатель & Значение \\
\hline
Количество документов & 12 \\
\hline
Успешно распарсено & 12 \\
\hline
Ошибок парсинга & 0 \\
\hline
Общий размер сырых HTML-документов & 2\,612\,218 байт \\
\hline
Общий размер выделенного текста & 32\,326 байт \\
\hline
Средний размер сырого документа & 217\,684.83 байт \\
\hline
Средний размер выделенного текста & 2\,693.83 байт \\
\hline
Всего слов & 6\,001 \\
\hline
Среднее число слов в документе & 500.08 \\
\hline
Уникальных текстов & 12 \\
\hline
Дубликатов & 0 \\
\hline
\end{tabular}
\end{center}

\subsection{Статистика по источникам}

\begin{center}
\small
\begin{tabular}{|l|r|r|r|r|r|}
\hline
Источник & Док. & Raw, байт & Text, байт & Avg raw, байт & Avg text, байт \\
\hline
Genius & 4 & 1\,664\,702 & 11\,432 & 416\,175.50 & 2\,858.00 \\
\hline
AZLyrics & 4 & 334\,505 & 10\,419 & 83\,626.25 & 2\,604.75 \\
\hline
Musixmatch & 4 & 613\,011 & 10\,475 & 153\,252.75 & 2\,618.75 \\
\hline
\end{tabular}
\end{center}

\section{Существующие поисковые системы}

Для выбранного корпуса уже существуют поисковые системы, которые можно использовать как внешний ориентир:

\begin{itemize}
    \item встроенный поиск Genius;
    \item встроенный поиск AZLyrics;
    \item встроенный поиск Musixmatch;
    \item Google или Яндекс с ограничением на сайт, например \code{site:genius.com Rakim It's Been a Long Time}.
\end{itemize}

Наличие таких поисковых систем подтверждает, что выбранный корпус подходит для дальнейших лабораторных работ: по документам уже можно выполнять поиск, а значит есть возможность сравнивать собственную систему с существующими решениями.

\subsection{Примеры запросов}

\begin{center}
\begin{tabular}{|p{0.28\linewidth}|p{0.27\linewidth}|p{0.35\linewidth}|}
\hline
Запрос & Где проверялось & Наблюдение \\
\hline
\code{Rakim its been a long time} & Musixmatch & Нужный трек не находится или находится хуже, чем при запросе с апострофом. В оригинальном названии используется форма \code{It's Been A Long Time}. Поиск чувствителен к пунктуации и хуже нормализует апострофы. \\
\hline
\code{Rakim It's Been a Long Time} & Musixmatch & При использовании апострофа запрос соответствует оригинальному названию трека, поэтому нужный документ находится корректнее. \\
\hline
\code{site:genius.com Rakim It's Been a Long Time} & Google/Яндекс & Поиск с ограничением на сайт позволяет найти страницу Genius напрямую и полезен как внешний способ проверки наличия документа в веб-индексе. \\
\hline
\code{Day N Nite Kid Cudi lyrics} & Встроенные поиски источников и Google/Яндекс & Запрос по исполнителю и названию обычно устойчивее, чем запрос только по фрагменту названия. \\
\hline
\end{tabular}
\end{center}

\placeholderfigure{Пример поиска Rakim без апострофа в Musixmatch}{
Здесь можно вставить скриншот: запрос \code{Rakim its been a long time}.\\
На скриншоте следует показать, что нужный трек находится хуже или отсутствует.
}

\placeholderfigure{Пример поиска Rakim с апострофом в Musixmatch}{
Здесь можно вставить скриншот: запрос \code{Rakim It's Been a Long Time}.\\
На скриншоте следует показать, что нужный трек находится корректнее.
}

\placeholderfigure{Пример поиска с ограничением site:}{
Здесь можно вставить скриншот Google или Яндекса с запросом\\
\code{site:genius.com Rakim It's Been a Long Time}.
}

\subsection{Недостатки существующего поиска}

По результатам ручной проверки можно выделить следующие недостатки:

\begin{itemize}
    \item поиск Musixmatch чувствителен к пунктуации: отсутствие апострофа в \code{It's} ухудшает выдачу по песне Rakim;
    \item разные источники по-разному записывают исполнителей с участием других артистов: \code{feat.}, \code{Ft.}, перечисление через запятую или амперсанд;
    \item точное совпадение по названию не всегда совпадает с пользовательским запросом: пользователи часто опускают апострофы, дефисы, скобки и дополнительные уточнения в названии;
    \item поиск по очень распространённым словам из текста песни может возвращать слишком много нерелевантных результатов;
    \item встроенные поиски сайтов зависят от их внутренней нормализации и ранжирования, поэтому одинаковый запрос может работать по-разному на Genius, AZLyrics и Musixmatch.
\end{itemize}

\section{Вывод}

В ходе лабораторной работы была подготовлена диагностическая выборка корпуса текстов песен из трёх источников: Genius, AZLyrics и Musixmatch. Для каждого источника была изучена структура HTML и способ хранения полезного текста. Реализовано извлечение текстов и метаданных, сохранение сырых HTML-документов и результатов разбора в JSON.

В выборке собрано 12 документов, все они успешно распарсены. Общий размер сырых HTML-документов составил 2\,612\,218 байт, общий размер выделенного текста --- 32\,326 байт. Средний размер сырого документа равен 217\,684.83 байт, средний размер выделенного текста --- 2\,693.83 байт.

Также были рассмотрены существующие поисковые системы по выбранным документам. На примере Musixmatch обнаружен недостаток поиска: запрос без апострофа \code{its been} хуже сопоставляется с оригинальным названием \code{It's Been}, что показывает важность нормализации пунктуации в будущей собственной поисковой системе.

\end{document}
